\chapter{Chapter09 — Continuous Reasonableness}
%%% generated by the blueprint pipeline from FABL.Chapter09.ContinuousReasonableness
\section{Overview}

Declarations of the Lean module \texttt{FABL.Chapter09.ContinuousReasonableness}.

\section{Declarations}

\begin{definition}[Fourth-moment measure condition]
\lean{FABL.IsReasonableMeasure}
\statementsource{Analysis-of-Boolean-Functions-by-Ryan-ODonnell}{pdf-ffbe9c7b3306-p248-b005}
\label{FABL.IsReasonableMeasure}
\leanok
Let $\Omega$ be a measurable space, let $\mu$ be a measure on it, let $B\in\mathbb{R}$,
and let $X:\Omega\to\mathbb{R}$. The measure and function satisfy the stated
fourth-moment condition when $\int X^4\,d\mu\le B\left(\int X^2\,d\mu\right)^2$.
\end{definition}

\begin{theorem}[Standard Gaussian second moment]
\lean{FABL.standardGaussian_secondMoment}
\label{FABL.standardGaussian_secondMoment}
\leanok
\difficulty{3}
The second moment of the standard Gaussian distribution is $1$, that is, $\int_{\mathbb
R}x^2\,d\gamma(x)=1$, where $\gamma$ is the standard Gaussian measure.
\end{theorem}

\begin{theorem}[Derivative of an exponential quadratic]
\lean{FABL.deriv_exp_half_sq}
\label{FABL.deriv_exp_half_sq}
\leanok
\difficulty{3}
For every real $t$, the derivative of $t\mapsto \exp(t^2/2)$ is $t\exp(t^2/2)$.
\end{theorem}

\begin{theorem}[Derivative of a Gaussian exponential function]
\lean{FABL.deriv_t_mul_exp_half_sq}
\label{FABL.deriv_t_mul_exp_half_sq}
\leanok
\uses{FABL.deriv_exp_half_sq}
\difficulty{3}
For every real $t$, the derivative of $t e^{t^2/2}$ is $(1+t^2)e^{t^2/2}$.
\end{theorem}

\begin{theorem}[Derivative of a Gaussian exponential polynomial]
\lean{FABL.deriv_one_add_sq_mul_exp_half_sq}
\label{FABL.deriv_one_add_sq_mul_exp_half_sq}
\leanok
\uses{FABL.deriv_exp_half_sq}
\difficulty{4}
For every real $t$, the derivative of $(1+t^2)e^{t^2/2}$ is $(3t+t^3)e^{t^2/2}$.
\end{theorem}

\begin{theorem}[Derivative of a cubic Gaussian exponential function]
\lean{FABL.deriv_three_mul_add_cube_mul_exp_half_sq}
\label{FABL.deriv_three_mul_add_cube_mul_exp_half_sq}
\leanok
\uses{FABL.deriv_exp_half_sq}
\difficulty{4}
For every real $t$, the derivative of $(3t+t^3)e^{t^2/2}$ is $(3+6t^2+t^4)e^{t^2/2}$.
\end{theorem}

\begin{theorem}[Standard Gaussian fourth moment]
\lean{FABL.standardGaussian_fourthMoment}
\label{FABL.standardGaussian_fourthMoment}
\leanok
\uses{FABL.deriv_exp_half_sq, FABL.deriv_one_add_sq_mul_exp_half_sq, FABL.deriv_t_mul_exp_half_sq, FABL.deriv_three_mul_add_cube_mul_exp_half_sq}
\difficulty{4}
The fourth moment of the standard Gaussian distribution is $3$, that is, $\int_{\mathbb
R} x^4\,d\gamma(x)=3$, where $\gamma$ is the standard Gaussian measure.
\end{theorem}

\begin{theorem}[A moment bound for the standard Gaussian]
\lean{FABL.standardGaussian_isThreeReasonable}
\proofsource{Analysis-of-Boolean-Functions-by-Ryan-ODonnell}{
  pdf-ffbe9c7b3306-p248-b005,
  pdf-ffbe9c7b3306-p248-b007
}
\label{FABL.standardGaussian_isThreeReasonable}
\leanok
\uses{FABL.IsReasonableMeasure, FABL.standardGaussian_fourthMoment, FABL.standardGaussian_secondMoment}
\difficulty{3}
For the standard Gaussian measure $\gamma$ on $\mathbb R$, its fourth moment is at most
three times the square of its second moment: $\int x^4\,d\gamma(x)\le 3\left(\int
x^2\,d\gamma(x)\right)^2$.
\end{theorem}

\begin{definition}[Uniform interval moment]
\lean{FABL.uniformIntervalMoment}
\label{FABL.uniformIntervalMoment}
\leanok
For a natural number $m$, define the $m$th moment of the uniform distribution on
$[-1,1]$ by $\frac12\int_{-1}^{1}x^m\,dx$.
\end{definition}

\begin{definition}[Uniform interval measure]
\lean{FABL.uniformIntervalMeasure}
\label{FABL.uniformIntervalMeasure}
\leanok
Define the uniform probability measure on $[-1,1]$ as the restriction of Lebesgue
measure to that interval, normalized to have total mass one.
\end{definition}

\begin{theorem}[Moments of the uniform interval measure]
\lean{FABL.integral_pow_uniformIntervalMeasure}
\label{FABL.integral_pow_uniformIntervalMeasure}
\leanok
\uses{FABL.uniformIntervalMeasure, FABL.uniformIntervalMoment}
\difficulty{4}
Let $X$ have the probability distribution obtained by restricting Lebesgue measure to
$[-1,1]$ and normalizing it. For every natural number $m$, its $m$th moment is
\[\mathbb E[X^m]=\frac12\int_{-1}^{1}x^m\,dx.\]
\end{theorem}

\begin{theorem}[Second moment of the uniform interval]
\lean{FABL.uniformIntervalMoment_two}
\label{FABL.uniformIntervalMoment_two}
\leanok
\uses{FABL.uniformIntervalMoment}
\difficulty{3}
The second moment of the uniform distribution on $[-1,1]$, defined by
$\frac12\int_{-1}^{1}x^2\,dx$, is $\frac13$.
\end{theorem}

\begin{theorem}[Fourth moment of the uniform interval]
\lean{FABL.uniformIntervalMoment_four}
\label{FABL.uniformIntervalMoment_four}
\leanok
\uses{FABL.uniformIntervalMoment}
\difficulty{3}
The fourth moment of the uniform distribution on $[-1,1]$, defined by
$\frac12\int_{-1}^{1}x^4\,dx$, is $\frac15$.
\end{theorem}

\begin{theorem}[Moment inequality for the uniform interval]
\lean{FABL.uniformInterval_isNineFifthsReasonable}
\proofsource{Analysis-of-Boolean-Functions-by-Ryan-ODonnell}{
  pdf-ffbe9c7b3306-p248-b005,
  pdf-ffbe9c7b3306-p248-b007
}
\label{FABL.uniformInterval_isNineFifthsReasonable}
\leanok
\uses{FABL.uniformIntervalMoment, FABL.uniformIntervalMoment_four, FABL.uniformIntervalMoment_two}
\difficulty{3}
For the uniform distribution on $[-1,1]$, the fourth moment is at most $\frac95$ times
the square of the second moment:
$\frac12\int_{-1}^{1}x^4\,dx\leq\frac95\left(\frac12\int_{-1}^{1}x^2\,dx\right)^2$.
\end{theorem}

\begin{theorem}[Fourth-moment bound for the uniform interval measure]
\lean{FABL.uniformIntervalMeasure_isNineFifthsReasonable}
\proofsource{Analysis-of-Boolean-Functions-by-Ryan-ODonnell}{
  pdf-ffbe9c7b3306-p248-b005,
  pdf-ffbe9c7b3306-p248-b007
}
\label{FABL.uniformIntervalMeasure_isNineFifthsReasonable}
\leanok
\uses{FABL.IsReasonableMeasure, FABL.integral_pow_uniformIntervalMeasure, FABL.uniformIntervalMeasure, FABL.uniformInterval_isNineFifthsReasonable}
\difficulty{3}
Let $\mu$ be Lebesgue measure conditioned on $[-1,1]$. For the identity random variable
$X(x)=x$, its fourth moment under $\mu$ is at most $\frac95$ times the square of its
second moment: $\int x^4\,d\mu(x)\leq\frac95\left(\int x^2\,d\mu(x)\right)^2$.
\end{theorem}
